\documentclass[12pt,a4paper]{article}
\usepackage{ugmath}
\usepackage{float}
\usepackage{placeins}
\usepackage{array}
\usepackage[labelfont=bf,labelsep=space]{caption}
\newcommand{\paren}[1]{\left( #1 \right)}
\newcommand{\alumno}{Ricardo León Martínez}
\newcommand{\materia}{Algebra Lineal I}
\newcommand{\profesor}{Claudia Reynoso Alcántara}
\newcommand{\tarea}{Tarea 5}
\newcommand{\fecha}{2/3/2026}

\begin{document}

    \begin{center}
        {\large\textbf{UNIVERSIDAD DE GUANAJUATO}}\\[0.3cm]
        {\normalsize\textbf{DIVISIÓN DE CIENCIAS NATURALES Y EXACTAS}}\\
        {\normalsize\textbf{CAMPUS GUANAJUATO}}\\[1cm]

        {\Large\textbf{\tarea\ (\materia)}}\\[1cm]
    \end{center}

    \noindent
    \textbf{Nombre:} \alumno 
    \hfill 
    \textbf{Fecha:} \fecha 
    \hfill 
    \textbf{Calificación:} \rule{3cm}{0.4pt}

    \vspace{0.3cm}

    Sea $F$ el campo $\mathbb{Q},\mathbb{R}$ o $\mathbb{C}$.

    \begin{mdframed}[style=mdbluebox,frametitle={Ejercicio 1}]
        Sea $V$ un espacio vectorial sobre $F$. Demuestra que
        \begin{enumerate}
            \item[(a)] El vector cero $0\in V$ es unico.
            \item[(b)] Para todos $v_{1},v_{2}\in V$, existe un único $v_{3}\in V$ tal que
            $v_{1}+v_{3}=v_{2}$.
            \item[(c)] Para todo $v\in V$, se tiene que $(-1)v=-v$ y $0v=0$. 
        \end{enumerate}
    \end{mdframed}

    \begin{proof}
        \begin{enumerate}
            \item[(a)] El vector cero $0\in V$ es unico.\\
            Sabemos por axiomas que existe un elemento $0\in V$ tal que
            \begin{equation*}
                v+0=v, \quad \forall v\in V.
            \end{equation*}
            Supongamos que existe $b\in V$ que tambien satisfacen la propiedad del neutro aditivo, es decir,
            \begin{equation*}
                v+b=v,\quad \forall v\in V.
            \end{equation*}
            Sumando $(-v)$ a ambos lados
            \begin{equation*}
                (-v)+(v+b)=(-v)+v.
            \end{equation*}
            Por asociatividad
            \begin{equation*}
                ((-v)+v)+b=0.
            \end{equation*}
            Luego
            \begin{equation*}
                0+b=0.
            \end{equation*}
            Por propiedad del elemento neutro
            \begin{equation*}
                b=0.
            \end{equation*}
            Así, el elemento neutro es único.
            \item[(b)] Para todos $v_{1},v_{2}\in V$, existe un único $v_{3}\in V$ tal que
            $v_{1}+v_{3}=v_{2}$.\\
            Sean $v_{1},v_{2}\in V$. Sabemos por axiomas que existe el inverso aditivo de $v_{1}$, denotado por
            $-v_{1}$, tal que
            \begin{equation*}
                v_{1}+(-v_{1})=0.
            \end{equation*}
            Tomemos
            \begin{equation*}
                v_{3}=(-v_{1})+v_{2}.
            \end{equation*}
            Entonces
            \begin{equation*}
                v_{1}+v_{3}=v_{1}+((-v_{1})+v_{2}).
            \end{equation*}
            Por asociatividad y definicion de inverso
            \begin{align*}
                v_{1}+((-v_{1})+v_{2})&=((-v_{1})+v_{1})+v_{2}\\
                &=0+v_{2}.
            \end{align*}
            Y por definicion de neutro aditivo
            \begin{equation*}
                0+v_{2}=v_{2}
            \end{equation*}
            Así, $v_{3}$ existe. Ahora supongamos que $w\in V$ satisface que
            \begin{equation*}
                v_{1}+w=v_{2}.
            \end{equation*}
            Entonces,
            \begin{equation*}
                v_{1}+v_{3}=v_{1}+w.
            \end{equation*}
            Sumamos $(-v_{1})$ a ambos lados
            \begin{equation*}
                (-v_{1})+(v_{1}+v_{3})=(-v_{1})+(v_{1}+w).
            \end{equation*}
            Por asociatividad
            \begin{equation*}
                ((-v_{1})+v_{1})+v_{3}=((-v_{1})+v_{1})+w.
            \end{equation*}
            Por propiedad del inverso aditivo
            \begin{equation*}
                0+v_{3}=0+w.
            \end{equation*}
            Por propiedad del neutro aditivo
            \begin{equation*}
                v_{3}=w.
            \end{equation*}
            Así, $v_{3}$ es unico.
            \item[(c)] Para todo $v\in V$, se tiene que $(-1)v=-v$ y $0v=0$.\\
            Notemos que
            \begin{equation*}
                1+(-1)=0.
            \end{equation*}
            Sea $v\in V$. Multplicamos $v$ a ambos lados
            \begin{equation*}
                (1+(-1))v=0v.
            \end{equation*}
            Por distributividad
            \begin{equation*}
                1v+(-1)v=0v.
            \end{equation*}
            En consecuencia
            \begin{equation*}
                v+(-1)v=0.
            \end{equation*}
            Luego,
            \begin{equation*}
                (-1)v=-v.
            \end{equation*}
            Ahora observamos que
            \begin{equation*}
                0=0+0.
            \end{equation*}
            Sea $v\in V$. Multplicamos $v$ a ambos lados
            \begin{equation*}
                0v=(0+0)v.
            \end{equation*}
            Por distributividad
            \begin{equation*}
                0v=0v+0v.
            \end{equation*}
            Sumamos $(-0v)$ a ambos lados. Y de esta forma por propiedad del inverso aditivo
            \begin{align*}
                0v&=0v-0v\\
                &=0
            \end{align*}
            Así, $(-1)v=-v$ y $0v=0$.
        \end{enumerate}
    \end{proof}

    \begin{mdframed}[style=mdbluebox,frametitle={Ejercicio 2}]
        Considera $V=\left\{(x_{1},x_{2}):x_{1},x_{2}\in\mathbb{C}\right\}$ y define las operaciones:
        \begin{align*}
            +:V\times V&\to V\\
            \left((x_{1},x_{2}),(y_{1},y_{2})\right)&\mapsto(x_{1}+y_{1}+1,x_{2}+y_{2}+1)
        \end{align*}
        \begin{align*}
            \cdot :\mathbb{C}\times V&\to V\\
            \left(c,(x_{1},x_{2})\right)&\mapsto(cx_{1}+c-1,cx_{2}+c-1).
        \end{align*}
        ¿Es $V$ un espacio vectorial sobre $\mathbb{C}$ con estas operaciones? (justifica tu respuesta).
    \end{mdframed}

    \begin{proof}
        Verifiquemos que $V$ satisface los axiomas de espacio vectorial sobre $\mathbb{C}$.\\
        \textbf{Cerrado bajo suma}.\\
        Si $u=(x_{1},x_{2}),v=(y_{1},y_{2})\in V$, entonces
        \begin{equation*}
            u+v=(x_{1}+y_{1}+1,x_{2}+y_{2}+1)\in V,
        \end{equation*}
        Pues $\mathbb{C}$ es cerrado bajo suma.\\
        \textbf{Asociatividad}\\
        Sea $u,v,w\in V$. Entonces
        \begin{align*}
            (u+v)+w&=(x_{1}+y_{1}+1+z_{1}+1,x_{2}+y_{2}+1+z_{1}+1),\\
            &=(x_{1}+y_{1}+z_{1}+2,x_{2}+y_{2}+z_{2}+2).
        \end{align*}
        Por otro lado,
        \begin{equation*}
            u+(v+w)=(x_{1}+y_{1}+z_{1}+2,x_{2}+y_{2}+z_{2}+2).
        \end{equation*}
        Luego, la suma es asociativa.\\
        \textbf{Conmutatividad}\\
        \begin{equation*}
            u+v=(x_{1}+y_{1}+1,x_{2}+y_{2}+1)=(y_{1}+x_{1}+1,y_{2}+x_{2}+1)=v+u.
        \end{equation*}
        \textbf{Existencia del neutro aditivo}\\
        Buscamos $(a,b)$ tal que
        \begin{equation*}
            (x_{1},x_{2})+(a,b)=(x_{1},x_{2}).
        \end{equation*}
        Esto exige que
        \begin{equation*}
            x_{1}+a=x_{1}, \quad x_{2}+b=x_{2}.
        \end{equation*}
        de donde $a=b=-1$. Luego el neutro es
        \begin{equation*}
            0_{V}=(-1,-1).
        \end{equation*}
        \textbf{Existencia de inverso aditivo}\\
        Dado $v=(x_{1},x_{2})$, buscmos $w=(a,b)$ tal que
        \begin{equation*}
            u+v=(-1,-1).
        \end{equation*}
        Esto exige que
        \begin{equation*}
            x_{1}+a+1=-1,\quad x_{2}+b+1=-1,
        \end{equation*}
        de donde
        \begin{equation*}
            a=-x_{1}-2,\quad b=-x_{2}-2.
        \end{equation*}
        Así, todo elemento tiene inverso.\\
        \textbf{Cerrado bajo el producto por escalares}\\
        Si $c\in\mathbb{C}$ y $c\in V$,
        \begin{equation*}
            cv=(cx_{1}+c-1,cx_{2}+c-1)\in V,
        \end{equation*}
        porque $\mathbb{C}$ es cerrado bajo suma y producto.\\
        \textbf{Distributividad respecto a suma en $V$}
        Sea $c\in\mathbb{C}$ y $u,v\in V$. Entonces
        \begin{align*}
            c(u+v)&=c((x_{1}+y_{1}+1,x_{2}+y_{2}+1)),\\
            &=(c(x_{1}+y_{1}+1)+c-1,c(x_{2}+y_{2}+1)+c-1),\\
            &=(cx_{1}+cy_{1}+2c-1,cx_{2}+cy_{2}+2c-1).
        \end{align*}
        Por otro lado,
        \begin{equation*}
            cu+cv=(cx_{1}+cy_{1}+2c-1,cx_{2}+cy_{2}+2c-1)
        \end{equation*}
        Así, el producto de es distributivo respecto a la suma en $V$.\\
        \textbf{Distributividad respecto a suma en $\mathbb{C}$}\\
        Sea $c,d\in\mathbb{C}$ y $v\in V$. Entonces
        \begin{equation*}
            (c+d)v=((c+d)x_{1}+(c+d)-1,(c+d)x_{2}+(c+d)-1).
        \end{equation*}
        Por otro lado,
        \begin{align*}
            cv+vd&=(cx_{1}+c-1,cx_{2}+c-1)+(dx_{1}+d-1,dx_{2}+d-1)\\
            &=((cx_{1}+c-1)+(dx_{1}+d-1)+1,(cx_{2}+c-1)+(dx_{2}+d-1)+1)\\
            &=((c+d)x_{1}+(c+d)-1,(c+d)x_{2}+(c+d)-1).
        \end{align*}
        Luego, el producto es distributivo respecto a la suma en $\mathbb{C}$.\\
        \textbf{Asociatividad respecto a el producto en $\mathbb{C}$}\\
        Sea $c,d\in\mathbb{C}$ y $v\in V$. Entonces
        \begin{equation*}
            (cd)v=((cd)x_{1}+(cd)-1,(cd)x_{2}+(cd)-1)
        \end{equation*}
        Por otro lado,
        \begin{align*}
            c(dv)&=c(dx_{1}+d-1,dx_{2}+d-1)\\
            &=(c(dx_{1}+d-1)+c-1,c(dx_{2}+d-1)+c-1)\\
            &=(cd)x_{1}+(cd)-1,(cd)x_{2}+(cd)-1.
        \end{align*}
        En consecuencia, hay asociatividad respecto al producto en $\mathbb{C}$.\\
        \textbf{Existencia del neutro multiplicativo}\\
        \begin{equation*}
            1(x_{1},x_{2})=(1\cdot x_{1}+1-1,1\cdot x_{2}+1-1)=(x_{1},x_{2}).
        \end{equation*}
        Así, existe el neutro multiplicativo. Finalmente concluimos que $V$ es un
        espacio vectorial sobre $\mathbb{C}$.
    \end{proof}

    \begin{mdframed}[style=mdbluebox,frametitle={Ejercicio 3}]
        Considera el espacio usual $\mathbb{R}^{n}$ sobre $\mathbb{R}$, determina cuáles de los siguientes
        subconjuntos son subespacios vectoriales de $\mathbb{R}^{n}$:
        \begin{enumerate}
            \item[(a)] $W=\left\{(a_{1},\dots,a_{n})\in\mathbb{R}^{n}:a_{1}\geq0\right\}$,
            \item[(b)] $W=\left\{(a_{1},\dots,a_{n})\in\mathbb{R}^{n}:a_{1}+3a_{2}=a_{3}\right\}$,
            \item[(c)] $W=\left\{(a_{1},\dots,a_{n})\in\mathbb{R}^{n}:a_{1}=a_{2}\right\}$,
            \item[(d)] $W=\left\{(a_{1},\dots,a_{n})\in\mathbb{R}^{n}:a_{j}\in\mathbb{Q}\text{ para todo }1\leq j\leq n\right\}$,  
        \end{enumerate}
        justifica tu respuesta.
    \end{mdframed}

    \begin{solucion}
        \begin{enumerate}
            \item[(a)] $W=\left\{(a_{1},\dots,a_{n})\in\mathbb{R}^{n}:a_{1}\geq0\right\}$\\
            Veremos si es un subespacio. Primero el vector cero pertence a $W$, pues $0\geq0$.
            Ahora tomemos
            \begin{equation*}
                v=(1,0,\dots,0)\in W.
            \end{equation*}
            Si multiplicamos por -1
            \begin{equation*}
                (-1)v=(-1,0,\dots,0),
            \end{equation*}
            y su primero coordenada es negativa, luego $(-1)v\notin W$. Por lo tanto, no es subespacio.
            \item[(b)] $W=\left\{(a_{1},\dots,a_{n})\in\mathbb{R}^{n}:a_{1}+3a_{2}=a_{3}\right\}$\\
            Veremos si es un subespacio. Primero, el vector cero satisface
            \begin{equation*}
                0+3\cdot0=0,
            \end{equation*}
            luego $0\in W$. Sea $u,v\in W$. Entonces
            \begin{equation*}
                u_{1}+3u_{2}=u_{3},\qquad v_{1}+3v_{2}=v_{3}.
            \end{equation*}
            Para $u+v$
            \begin{equation*}
                (u_{1}+v_{1})+3(u_{2}+v_{2})=(u_{1}+3_{2})+(v_{1}+3v_{2})=u_{3}+v_{3}.
            \end{equation*}
            Luego $u+v\in W$. Sea $c\in\mathbb{R}$. Entonces
            \begin{equation*}
                cu_{1}+3cu_{2}=c(u_{1}+3u_{2})=cu_{3}.
            \end{equation*}
            Luego $cu\in W$. Por tanto, si es subespacio.
            \item[(c)] $W=\left\{(a_{1},\dots,a_{n})\in\mathbb{R}^{n}:a_{1}=a_{2}\right\}$\\
            Veamos si es un subespacio. Primero el vector cero cumple 0=0 entonces $0\in W$.
            Si $u,v\in W$, entonces
            \begin{equation*}
                u_{1}=u_{2},\qquad v_{1}=v_{2}.
            \end{equation*}
            Entonces
            \begin{equation*}
                (u_{1}+v_{1})=(u_{2}+v_{2}).
            \end{equation*}
            Luego, es cerrado bajo suma,
            \begin{equation*}
                cu_{1}=cu_{2}.
            \end{equation*}
            Luego es cerrado bajo escaleres. Por tanto si es un subespacio.
            \item[(d)] $W=\left\{(a_{1},\dots,a_{n})\in\mathbb{R}^{n}:a_{j}\in\mathbb{Q}\text{ para todo }1\leq j\leq n\right\}$\\
            Veamos si es un subespacio. El vector cero pertenece a $W$. Es cerrado bajo la suma pues la suma de racionales es racional.
            Ahora tomemos
            \begin{equation*}
                v=(1,0,\dots,0)\in W.
            \end{equation*}
            Multiplicando por $\sqrt{2}\in\mathbb{R}\setminus\mathbb{Q}$,
            \begin{equation*}
                \sqrt{2}v=(\sqrt{2},0,\dots,0),
            \end{equation*}
            que no pertenece a $W$. Por tanto no es un subespacio.
        \end{enumerate}
    \end{solucion}

    \begin{mdframed}[style=mdbluebox,frametitle={Ejercicio 4}]
        Sean $W_{1},W_{2}$ subespacios vectoriales de un espacio vectorial $V$ sobre el campo
        $F$. Demuestra que si $W_{1}\cup W_{2}$ es un subespacio vectorial $V$ entonces uno de los espacios
        $W_{j}$ está contenido en el otro.
    \end{mdframed}

    \begin{proof}
        Supongamos que $W_{1}\cup W_{2}$ es subespacio y que
        \begin{equation*}
            W_{1}\nsubseteq W_{2}\quad\text{ y }\quad W_{2}\nsubseteq W_{1}
        \end{equation*}
        Entonces existen
        \begin{equation*}
            u\in W_{1}\setminus W_{2},\qquad v\in W_{1}\setminus W_{2}.
        \end{equation*}
        Como $W_{1}\cup W_{2}$ es subespacio, es cerrado bajo suma, luego
        \begin{equation*}
            u+v\in W_{1}\cup W_{2}.
        \end{equation*}
        Por definición de unión, entonces ocurre una de dos cosas. Primer caso, $u+v\in W_{1}$.
        Como $u\in W_{1}$ y $W_{1}$ es subespacio, entonces también
        \begin{equation*}
            (u+v)-u\in W_{1}.
        \end{equation*}
        Pero
        \begin{equation*}
            (u+v)-u=v.
        \end{equation*}
        Luego $v\in W_{1}$, lo cual contradice que $v\notin W_{1}$. Segundo caso, $u+v\in W_{2}$.
        Como $v\in W_{2}$, entonces
        \begin{equation*}
            (u+v)-v\in W_{2}.
        \end{equation*}
        Pero
        \begin{equation*}
            (u+v)-u=u.
        \end{equation*}
        Luego $u\in W_{2}$, contradicción. En ambos casos llegamos a contradicción Luego
        \begin{equation*}
            W_{1}\subseteq W_{2}\quad\text{ o }\quad W_{2}\subseteq W_{1}
        \end{equation*}
    \end{proof}

    \begin{mdframed}[style=mdbluebox,frametitle={Ejercicio 5}]
        Considera el $\mathbb{R}$-espacio vectorial, $V=\left\{f:\mathbb{R}\to\mathbb{R}:f\text{ es una función continua}\right\}$.
        Demuestra que la función $f(x)=\sin(x)$ no es combinación lineal de las funciones $x,x^{2},x^{3}$ (puedes usar todo lo que
        sabes de tus clases de calculo)
    \end{mdframed}

    \begin{proof}
        SUpongamos que existen $a,b,c\in\mathbb{R}$ tales que
        \begin{equation*}
            \sin(x)=ax+bx^{2}+cx^{3}
        \end{equation*}
        sabemos que
        \begin{equation*}
            \sin(x)=\sum_{n=0}^{\infty}\frac{(-1)^{n}}{(2n+1)!}x^{2n+1}=x-\frac{x^{3}}{3!}+\frac{x^{5}}{5!}-\frac{x^{7}}{7!}+\cdots
        \end{equation*}
        Esta expansión converge para todo $x\in\mathbb{R}$. Si dos funciones analíticas coinciden en un intervalo,
        entonces coinciden sus series de taylor termino a termino. Por tanto si la igualdad fuese cierta, deberia cumplirse
        \begin{equation*}
            x-\frac{x^{3}}{6}+\frac{x^{5}}{120}-\cdots=ax+bx^{2}+cx^{3}.
        \end{equation*}
        Pero entonces en el término $x^{5}$ deberia tener coeficiente 0 en el polinomio,
        mientras que en la seria de $\sin(x)$ el coeficiente es
        \begin{equation*}
            \frac{1}{120}\neq0.
        \end{equation*}
        Contradicción. Por tanto no puede existir tal combinación lineal.
    \end{proof}

\end{document}